\begin{proposition}
\label{editorial-source-proposition-localize-twice-module}
View $S'^{-1}M$ as an $A$-module, then $S^{-1}(S'^{-1}M)$ is
isomorphic to $(SS')^{-1}M$.
\end{proposition}

\begin{proof}
Note that given an $A$-module $M$, the universal property of
Lemma \ref{algebra-lemma-universal-property-localization-module} is available.
We give an explicit construction of the isomorphism.

\medskip\noindent
We define the maps as follows
\begin{align*}
& f : S^{-1}(S'^{-1}M) \longrightarrow (SS')^{-1}M, \quad \frac{x/s'}{s}\mapsto
x/ss'\\
& g : (SS')^{-1}M \longrightarrow S^{-1}(S'^{-1}M), \quad x/t\mapsto
\frac{x/s'}{s}\ \text{for some }s\in S, s'\in S', \text{ and }
t = ss'
\end{align*}
We have to check that these homomorphisms are well-defined, that is,
independent of the choice of the fraction. The universal property proves that these maps are well defined:
every element of $S'$ and of $S$ acts invertibly on $(SS')^{-1}M$,
so the map $M \to (SS')^{-1}M$ extends first to $S'^{-1}M$
and then to $S^{-1}(S'^{-1}M)$, giving $f$ with the displayed formula.
On $S^{-1}(S'^{-1}M)$ every $s' \in S'$ acts invertibly already
on the inner localization, and every $s \in S$ acts invertibly
on the outer localization. Thus every $ss' \in SS'$ acts invertibly,
and $M \to S^{-1}(S'^{-1}M)$ extends uniquely to $g$.
The displayed formula for $g$ follows by inverting the actions of
$s'$ and $s$; it therefore does not depend on the chosen factorization
of the denominator. Both composites restrict to the original
localization map on $M$. Uniqueness in the universal property,
applied once for $(SS')^{-1}M$ and successively for the two
localizations, gives $f \circ g = \mathrm{id}_{(SS')^{-1}M}$ and
$g \circ f = \mathrm{id}_{S^{-1}(S'^{-1}M)}$.
\end{proof}